\chapter{Bilan, difficultés et perspectives}
\label{ch:bilan}

\section{Satisfactions}

Le projet a permis de construire et de \emph{posséder} une application web multicouche complète :
un socle « framework » écrit à la main (routeur, middlewares, logger PSR-3, utilitaires), une
sécurité pensée couche par couche (XSS, CSRF, injection SQL, IDOR), une base PostgreSQL conçue
proprement avec migrations idempotentes, et une chaîne Docker reproductible du développement à la
production. La migration de MariaDB vers PostgreSQL et l'audit d'architecture ont particulièrement
consolidé la compréhension du fonctionnement interne de l'application.

\section{Difficultés rencontrées et résolues}

Ces difficultés démontrent la \textbf{démarche structurée de résolution de problème} (compétence
transversale T2) : chacune suit le schéma symptôme $\to$ diagnostic $\to$ correctif.

\begin{description}
  \item[\emph{Headers already sent} à la connexion.] La réponse JSON était émise \emph{avant}
        \texttt{session\_start()}, empêchant la pose du cookie \texttt{PHPSESSID}.
        \emph{Correctif} : démarrer la session dans le \emph{bootstrap}
        (\file{public/index.php}), avant toute sortie possible.
  \item[Requête N+1 sur le fil.] L'affichage des auteurs déclenchait une requête par message.
        \emph{Correctif} : préchargement en une requête \texttt{IN} (section~\ref{sec:acces}),
        avec gestion du cas \texttt{IN ()} vide, interdit en PostgreSQL.
\end{description}

\section{Perspectives}

Les axes d'amélioration, déjà identifiés dans le dossier, sont :

\begin{itemize}
  \item le \textbf{durcissement de la validation serveur} (format de l'email, robustesse du mot de
        passe) et l'\textbf{anti-fixation de session} (\texttt{session\_regenerate\_id}) ;
  \item la \textbf{validation du téléversement par type MIME réel} (\texttt{finfo} /
        \texttt{getimagesize}) en complément de l'extension ;
  \item une passe d'\textbf{accessibilité} (RGAA) plus poussée ;
  \item la \textbf{maintenabilité du front-end} : séparation et factorisation des feuilles de
        style et des scripts, \textbf{chargement conditionnel du JavaScript} selon le gabarit
        (\emph{layout}) plutôt qu'un script global, et mise en place d'un \emph{lint} JavaScript
        (ESLint).
\end{itemize}

